\documentclass[11pt,a4paper]{article}
\usepackage[margin=25mm]{geometry}
\usepackage{amsmath,amssymb,amsthm}
\setlength{\emergencystretch}{2em}
\newtheorem{lemma}{Lemma}
\newtheorem{theorem}[lemma]{Theorem}
\newcommand{\gen}[1]{\langle#1\rangle}
\newcommand{\Alt}{\operatorname{Alt}}
\newcommand{\Sym}{\operatorname{Sym}}
\title{Pure alternating kernels in the critical prefixes}
\author{Human Compression project: L044 working continuation}
\date{}
\begin{document}\maketitle
We retain the full critical six-point-head hypotheses, actual ambient
IP and smaller-degree induction. Thus $n=2r\ge18$, the first head
is supported on $D=\{1,\ldots,6\}$, and
\[
 \tau=\rho_3=(6\ b)\delta,\quad \sigma=\rho_4,\quad
 R=G_{\ge3}=\Alt(Z),\qquad Z=(\Omega\setminus D)\cup\{6\}.
\]
The head is transitive on $D$, $\sigma$ fixes $D$, and $G_{\ge5}$
fixes $D\cup\{b\}$. The complete critical-tail proof gives
$x=b^\sigma\ne b$ and $\tau x=x$ without assuming $G_{\ge4}$
transitive. The involution $\delta$ is odd and fixes $D\cup\{b\}$.
Every alternating group below is the actual natural subgroup fixing
all points outside its displayed support.

\begin{lemma}[A support-merging extension]\label{merge}
Let $P$ preserve $C$, contain $\Alt(C)$, and have $|C|\ge4$.
Let $t$ be an involution with $|C\cap C^t|\ge2$.
Then $\gen{P,t}$ contains $\Alt(C^{\gen{P,t}})$, where the
superscript denotes the union of all translates of $C$.
\end{lemma}
\begin{proof}
Natural alternating groups on two supports of size at least four
whose intersection contains two points generate the alternating
group on their union. Indeed, choose two common points $a,b$.
The three-cycles $(a\ b\ z)$ for every other point in the union
belong to the generated group; these generate its full alternating
group. The latter fact follows by conjugating such cycles to obtain
all three-cycles, which generate every even permutation.

Starting with $C$, first adjoin $C^t$ using this observation.
Maintain a set $U$ containing $C\cup C^t$ with $\Alt(U)$ in
$\gen{P,t}$. For $p\in P$, both $U$ and $U^p$ contain $C$;
for $t$, both $U$ and $U^t$ contain $C\cup C^t$.
Thus we may adjoin the translate by any generator of $P$ or by
$t$ and still have the full alternating group on the enlarged set.
Finite iteration terminates at an invariant set, exactly the
union of all translates of $C$. This proves the assertion without
a primitive-group or Jordan theorem.
\end{proof}

\begin{theorem}[Actual prefix kernels and their exterior groups]\label{prefix}
Put $P_k=G_{\le k}$ and let $C_k$ be its actual orbit containing
$D$. Then
\[
 C_3=D\cup\{b\},\qquad C_4=D\cup\{b,x\},\qquad
               \Alt(C_k)\le P_k\quad(3\le k\le r-1).
\]
Let $J_k$ be the literal permutation image of $P_k$ on
$\Omega\setminus C_k$. With restrictions to these disjoint sets,
\[
 P_k=\{(s,j)\in\Sym(C_k)\times J_k:
                              \operatorname{sgn}s=\operatorname{sgn}j\}.
                                                        \tag{1}
\]
For $k\ge4$ every generator of index at least $k+2$ fixes $C_k$
pointwise. The sets $C_3,C_4,\ldots,C_{r-1}$ strictly increase,
and $C_{r-1}=\Omega$.
\end{theorem}
\begin{proof}
Conjugating the single core transposition $(6\ b)$ by the
transitive head on $D$ gives every $(d\ b)$ in the core image.
Products of two conjugates of $\tau$ cancel the exterior $\delta$
and give the pure three-cycles on $D\cup\{b\}$.
Thus $P_3$ preserves this seven-point set, is transitive there,
and contains its full alternating group. The pure-overlap/IP
argument gives $\tau x=x$, so every generator of $P_3$ fixes
$x$. The permutation $\sigma$ fixes $D$ and exchanges $b,x$.
Consequently $P_4$ preserves $D\cup\{b,x\}$ and contains its
alternating group, by merging the conjugate seven-point supports.
This proves both initial orbit identities.

Suppose $k\ge5$ and the conclusion is proved for $k-1$ and
$k-2$. String commutation makes $\rho_k$ centralize
$\Alt(C_{k-2})$, and hence fix that set pointwise: the centralizer
of natural $\Alt(C)$ on its support is trivial when $|C|\ge4$.
Here $|C_{k-2}|\ge7$, so $C_{k-1}$ and its $\rho_k$-translate
intersect in at least seven points. Lemma~\ref{merge}, applied
to $P_{k-1}$, gives $\Alt(C_k)\le P_k$. Its invariant union is
exactly the orbit $C_k$, since $P_{k-1}$ is transitive on $C_{k-1}$.
The same centralizer argument proves the stated fixation by every
later distant generator.

For (1), all elements of $P_k$ are even, so the two restriction
signs agree. Given any $j\in J_k$, choose one lift $(s,j)\in P_k$.
Multiplication by the full kernel $\Alt(C_k)\times1$ supplies
exactly all permutations on $C_k$ with sign equal to that of $j$.
This proves equality with the displayed set, including the case
that every element of $J_k$ is even.

If $C_k=C_{k-1}$ for some $k\ge4$, then $P_k$ preserves
$C_{k-1}$ and all generators of index at least $k+1$ fix it,
because they centralize $\Alt(C_{k-1})$. Transitivity of $G$
would force $C_{k-1}=\Omega$. But then $P_{k-1}$ contains
$\Alt(\Omega)$, equals $G$, and contradicts independence of the
actual string C-marking. Thus equality cannot occur. Finally
$C_{r-1}=\Omega$ by transitivity, proving all assertions.
\end{proof}

\begin{lemma}[The exact local omission bridge]\label{orbit}
For $j\ge3$, put $R_j=\gen{\rho_k:k\ge3,\ k\ne j}$ and
let $O_6$ be its orbit of $6$ on $Z$. The $G_j$-orbits are
$O_6\cup(D\setminus\{6\})$ and all other $R_j$-orbits on $Z$.
Consequently every $j$-edge has the same eligibility for its
ambient omission $G_j$ and its actual local omission $R_j$.
Moreover, for $k\ge3$,
\[
 G_{\{3,\ldots,k\}}=P_k\cap R
 =\{(s,j)\in\Sym(C_k\setminus(D\setminus\{6\}))\times J_k:
                                \operatorname{sgn}s=\operatorname{sgn}j\}.
                                                        \tag{2}
\]
\end{lemma}
\begin{proof}
We have $G_j=\gen{G_{<3},R_j}$. The head moves only $D$, meets
$Z$ at the single vertex $6$, and is transitive on $D$. It therefore
adjoins the other five head vertices to $O_6$ and merges no other
$R_j$-orbits. Since $\rho_j$ is supported on $Z$, eligibility is
unchanged. Formula (2) is (1) restricted to permutations fixing
$D\setminus\{6\}$ pointwise; $R$ is exactly that natural alternating
group, and actual ambient IP supplies the marked intersection.
No quotient intersection property is asserted.
\end{proof}

\paragraph{Attribution and remaining scope.}
Root proposed the initial eight-point parity-fiber description and
suggested extending the pure support. The finite support-merging
induction and complete orbit/fiber formulas above implement that
suggestion; a distinct review must disclose this contribution.
All complete critical-head and pure-overlap proofs remain paid.
At a later unshared cut with transitive supported prefix, (1) now
identifies that prefix as natural alternating directly, because its
exterior image is trivial. A shared cut may retain a nontrivial
$J_k$; it is not silently deleted. The index-$3$ prefix and the
remaining opposite-end and source rigidity obligations are not
resolved by these formulas. No new computation, formalization,
whole-original credit or question claim is made.
\end{document}
